\setcounter{section}{6}
\section{Hệ trục tọa độ trong không gian}
\subsection{Đề 1}
\Opensolutionfile{ans}[ans/ans-2-B7]
\caulc
\begin{ex}%[Câu 1]%[2H2N2-2]
	Trong không gian $Oxyz$, cho điểm $ M\left(1;-2;3\right) $. Chọn khẳng định đúng trong các khẳng định sau.
	\choice
	{\True $ \overrightarrow{OM}= \overrightarrow{i}- 2\overrightarrow{j}+ 3\overrightarrow{k}$}
	{$ \overrightarrow{MO}= \overrightarrow{i}- 2\overrightarrow{j}+ 3\overrightarrow{k}$}
	{$ \overrightarrow{OM}= \overrightarrow{k}- 2\overrightarrow{j}+ 3\overrightarrow{i}$}
	{$ \overrightarrow{OM}= \overrightarrow{j}- 2\overrightarrow{i}+ 3\overrightarrow{k}$}
	\loigiai{Theo định nghĩa ta có $ \overrightarrow{OM}= \overrightarrow{i}- 2\overrightarrow{j}+ 3\overrightarrow{k}$.
	}
\end{ex}

\begin{ex}%[Câu 2]%[2H2N2-3]
	Trong không gian $Oxyz$, tọa độ của véc-tơ $ \overrightarrow{i} $ là
	\choice
	{\True $ \left(1;0;0\right) $}
	{$ \left(0;1;0\right) $}
	{$ \left(0;0;1\right) $}
	{$ \left(0;1;1\right) $}
	\loigiai{
Véc-tơ 	$ \overrightarrow{i} $ là véc-tơ đơn vị của trục $ Ox $ nên $ \overrightarrow{i} =\left(1;0;0\right)$.
}
\end{ex}
\begin{ex}%[Câu 3]%[2H2N2-3]
	Trong không gian $Oxyz$, cho $\overrightarrow{a}= -\overrightarrow{i}+ 2\overrightarrow{j}- 3\overrightarrow{k}$. Tọa độ của véc-tơ $\vec{a}$ là
	\choice
	{$(2;-1;-3)$}
	{$(-3;2;-1)$}
	{$(2;-3;-1)$}
	{\True $(-1;2;-3)$}
	\loigiai{
		Ta có $\vec{a}= -\overrightarrow{i}+ 2\overrightarrow{j}- 3\overrightarrow{k} \Rightarrow \overrightarrow{a}= (-1;2;-3)$.
	}
\end{ex}
\begin{ex}%[Câu 4]%[2H2N2-3]
	Trong không gian $Oxy$, cho $A(1;-1;2)$ và $B(-1;0;1)$. Tọa độ véc-tơ $\overrightarrow{AB}$ là
	\choice
	{$ (2;-1;1) $}
	{$ (-2;-1;-1) $}
	{\True $ (-2;1;-1) $}
	{$ (0;-1;3) $}
	\loigiai{
		Ta có $\overrightarrow{AB}=\left(-1-1;0-(-1);1-2\right)=(-2;1;-1)$. 
	}
\end{ex}
\begin{ex}%[Câu 5]%[2H2H2-2]
	Trong không gian $Oxyz$, cho điểm $A(2;-1;3)$. Hình chiếu của $A$ trên trục $Oz$ là
	\choice
	{$Q(2;-1;0)$}
	{\True $P(0;0;3)$}
	{$N(0;-1;0)$}
	{$M(2;0;0)$}
	\loigiai{
		Hình chiếu của điểm $A(2;-1;3)$ trên trục $Oz$ có tọa độ là $(0;0;3)$.
	}
\end{ex}
\begin{ex}%[Câu 6]%[2H2H2-3]
	Trong không gian $Oxyz$, cho $\overrightarrow{OA}=(3;1;-2)$. Tìm tọa độ điểm $A$.
	\choice
	{\True$\left(3;1;-2\right)$}
	{$\left(-3;-1;2\right)$}
	{$\left(\dfrac{3}{2};\dfrac{1}{2};-1\right)$}
	{$\left(\dfrac{-3}{2};\dfrac{-1}{2};1\right)$}
	\loigiai{Theo định nghĩa tọa độ điểm $ A $ là $ \left(3;1;-2\right) $.
	}
\end{ex}
\begin{ex}%[Câu 7]%[2H2H2-2]
	Trong không gian $ Oxyz $, điểm nào sau đây thuộc trục $ Oy $?
	\choice
	{\True $ M\left(0;-\dfrac{1}{3};0\right) $}
	{$ N\left(0;0;-\dfrac{1}{3}\right) $}
	{$ P\left(-\dfrac{1}{3};0;0\right) $}
	{$ Q\left(-\dfrac{1}{3};0;\dfrac{1}{3}\right) $}
	\loigiai{Điểm thuộc trục $ Oy $ có tọa độ dạng $ (0;y;0) $ nên điểm $ M\left(0;-\dfrac{1}{3};0\right) $ thuộc $ Oy $.
	}
\end{ex}
\begin{ex}%[Câu 8]%[2H2H2-2]
	Trong không gian $Oxyz$, cho điểm $A\left(1;-3;1\right)$. Hình chiếu vuông góc của điểm $A$ trên mặt phẳng $\left(Oyz\right)$ là điểm
	\choice
	{$M\left(1;0;0\right)$}
	{\True $N\left(0;-3;1\right)$}
	{$P\left(0;-3;0\right)$}
	{$Q\left(0;0;1\right)$}
	\loigiai{
		Khi chiếu vuông góc một điểm trong không gian lên mặt phẳng $(Oyz)$, ta giữ lại các thành phần tung độ và cao độ nên hình chiếu của điểm $A(1;-3;1)$ lên $(Oyz)$ là điểm $N(0;-3;1)$.	
	}
	
\end{ex}
\begin{ex}%[Câu 9]%[2H2H2-3]
	Trong không gian $Oxyz$, cho điểm $N(3;1;0)$ và $\overrightarrow{MN}=(-1;-1;0)$. Tìm tọa độ của điểm $M$.
	\choice
	{$M(4;0;0)$}
	{$M(-4;-2;0)$}
	{\True $M(4;2;0)$}
	{ $M(2;0;0)$}
	\loigiai{
		Từ $N(3;1;0)$ và $\overrightarrow{MN}=(-1;-1;0)$ ta có
		$\heva{&3-x_M=-1\\&1-y_M=-1\\&0-z_M=0}\Leftrightarrow\heva{&x_M=4\\&y_M=2\\&z_M=0}$. 
		Vậy $M(4;2;0)$.
	}
\end{ex}
\begin{ex}%[Câu 10]%[2H2H2-2]
Trong không gian $Oxyz$, tọa độ điểm $ A $ thuộc tia $ Ox $ và độ dài đoạn $ OA=3 $ là
\choice
{\True $ (3;0;0) $}
{ $ (-3;0;0) $}
{ $ (0;3;0) $}
{$ (1;1;1) $}
\loigiai{Vì $ A $ thuộc tia $ Ox $ nên $ \overrightarrow{OA}=3\overrightarrow{i} $.\\
Vậy $ A\left(3;0;0\right) $.
}
\end{ex}

\begin{ex}%[Câu 11]%[2H2V2-2]
\immini{Cho hình lập phương $A B C D . A' B' C' D'$. Có thể lập một hệ toạ độ $Oxyz$ có gốc $ O $ trùng với đỉnh $C$ và $ B $, $ D $, $ C'$ lần lượt thuộc trục $ Ox $, $ Oy $, $ Oz $. Giả sử $ A'(3 ; 2 ;4)$ . Tìm tọa độ đỉnh $ B' $.}
	{\begin{tikzpicture}[scale=0.6, font=\footnotesize, line join=round, line cap=round, >=stealth]
			\def\bc{4} % cạnh BC
			\def\ba{2} % cạnh BA
			\def\gocB{35} % góc B của đáy
			\coordinate[label=below right:$B$] (B) at (0,0);
			\coordinate(C) at (\gocB:\ba);
			\coordinate[label=below:$A$] (A) at (\bc,0);
			\coordinate[label=below right:$D$] (D) at ($(C)-(B)+(A)$);
			\coordinate[label=above left:$A'$] (A') at ($(A)+(90:\bc)$);
			\coordinate[label=left:$B'$] (B') at ($(B)-(A)+(A')$);
			\coordinate[label=below right:$C'$] (C') at ($(C)-(A)+(A')$);
			\coordinate[label=right:$D'$] (D') at ($(D)-(A)+(A')$);
			\draw (B')--(B)(D)--(D')--(A')--(B')--(C')--(D') (B)--(A)--(D) (A)--(A');
			\draw[dashed] (C)--(B)(C)--(D) (C)--(C');
			\foreach \diem in {A,B,C,D,A',B',C',D'}	\fill (\diem)circle(1.0pt);
			\draw[-stealth] (D)--+(0:1)node[right]{$y$};
			\draw[-stealth] (C')--+(90:0.8)node[above]{$z$};
			\draw[-stealth] (B)--+(-145:.75)node[below]{$x$};
			%\draw[-stealth,blue] (C)--+(0:.8)node[above]{$\vec{j}$};
			%\draw[-stealth,blue] (C)--+(90:.85)node[above left]{$\vec{k}$};
			%\draw[-stealth,blue] (C)--+(-145:.65)node[above]{$\vec{i}$};
			\fill(C)circle(1pt) node[below right]{$C\equiv O$};			
	\end{tikzpicture}	}
\choice
{$B'(3;-2;4)$}
{$B'(3;2;0)$}
{$B'(0;2;4)$}
{\True $B'(3;0;4)$}
\loigiai{
Dễ thấy $ B' $ là hình chiếu của $ A' $ lên mặt phẳng tọa độ $ \left(Oxz\right) $ nên $B'(3;0;4)$.
}
\end{ex}
\begin{ex}%[Câu 12]%[2H2V2-2]
\immini{Trong không gian $Oxyz$, tìm tọa độ của các vectơ $\overrightarrow{A _1A}$ ở hình vẽ sau đây, biết $ A\left(4;5;3\right) $.}
	{\begin{tikzpicture}[scale=.5]
			\def\goc{215}\def\a{4}\def\b{5}\def\c{3}
			\path
			(0:0) coordinate (O)
			(0:\b) coordinate (K)
			++(\goc:\a) coordinate (A_1)
			++(90:\c) coordinate (M)
			(\goc:\a) coordinate (H)
			(K)++(90:\c) coordinate (A_2)
			(90:\c) coordinate (L)
			++(\goc:\a) coordinate (A_3)
			;
			\foreach \x in {1,...,9}\draw[gray,opacity=.5] 
			(0:\x) --++ (90:9) (90:\x)--++(0:9)
			(0,\x)--++(\goc:9) (\goc:\x)--++(90:9)
			(\goc:\x)--++(0:9) (0:\x)--++(\goc:9)
			;
			\draw[-stealth] (O)--++(\goc:7) node[left]{$x$};
			\draw[-stealth] (O)--++(0:7) node[below]{$y$};
			\draw[-stealth] (O)--++(90:5) node[left]{$z$};
			\draw 
			(H)--(A_1)--(K)--(A_2)--(L)--(A_3)--cycle
			(A_3)--(M)--(A_2) (M)--(A_1)
			;
			%\draw[red] (L)--(M)--(H) (M)--(K);
			\draw[dashed] (O)--(H);
			\draw [dashed](O)--(L);
			\draw (A_2)--(M);
			\draw (A_1)--(M);
			\path
			node at ($(M)+(1.5,-0.2)$)[left]{$A$}
			node at (K)[above right]{$5$}
			node at (H)[below]{$4$}
			node at (L)[above right]{$3$}
			node at (A_3)[above left]{$A_3$};
			\foreach \x/\g in {O/180,H/180,A_1/270,K/-60,L/135,A_2/60}\draw[fill=white] (\x) circle (.05) +(\g:.5) node{$\x$};
	\end{tikzpicture}}
\choice
{$ (0,0,3) $}
{\True $ (0,0,-3) $}
{$ (3,0,0) $}
{$ (0,3,0) $}
\loigiai{
Ta có $\overrightarrow{AA_1} =\overrightarrow{LO}=\left(0;0;-3\right)$.}
\end{ex}
\Closesolutionfile{ans}

\inputansbox[3]{6}{ans/ans-2-B7}

\cauds
\Opensolutionfile{ans}[ans/ans-2-B7-DS]
\setcounter{ex}{0}
\begin{ex}%[Câu 1]%[2H2H2-3]
Trong không gian $Oxyz$, cho hai điểm $B\left( 0;3;1\right)$, $C\left( -3;6;4\right)$. Các khẳng định sau đúng hay sai?
\choiceTF
{Véc-tơ đơn vị của trục $ Oy $ là $ \overrightarrow{j}=(0;-1;0) $}
{$ \overrightarrow{BC}=\left(3;-3;-3\right) $}
{\True Hình chiếu của $ C $ trên $ \left(Oxy\right) $ có tọa độ là $ \left(-3;6;0\right) $}
{\True Gọi $M$ là điểm nằm trên đoạn $BC$ sao cho $MC=2MB$. Tọa độ của $ M $ là $ \left( -1;4;2\right) $}
\loigiai{
	\begin{itemchoice}
	\itemch Véc-tơ đơn vị của trục $ Oy $ là $ \overrightarrow{j}=(0;1;0) $.
	\itemch $ \overrightarrow{BC}=\left(x_C-x_B;y_C-y_B;z_C-z_B\right)\left(-3;3;3\right) $.
	\itemch Hình chiếu của $ C $ trên $ \left(Oxy\right) $ có tọa độ là $ \left(-3;6;0\right) $
	\itemch Gọi $M(a;b;c)$, khi đó $\overrightarrow{MC}=-2\overrightarrow{MB} \Leftrightarrow \heva{&-3-a=-2(-a)\\ &6-b=-2(3-b)\\ &4-c=-2(1-c)} \Leftrightarrow \heva{&a=-1\\ &b=4\\ &c=2.}$\\
	Vậy $M\left( -1;4;2\right)$.
	\end{itemchoice}
}
\end{ex}
\begin{ex}%%[Câu 2]%[2H2H2-3]
	Trong không gian $Oxyz$, cho véc-tơ $\overrightarrow{OA}=2\vec{i}+5\vec{j}+3\vec{k}$. Các khẳng định sau đúng hay sai?
	\choiceTF
	{\True Tọa độ điểm $A$ là $(2;5;3)$}
	{Hình chiếu vuống góc của $A$ lên trục $Ox$ là $A'(0;5;0)$}
	{Hình chiếu vuông góc của $A$ lên mặt phẳng $(Oyz)$ là $H(2;0;0)$}
	{\True Điểm $ A' $ đối xứng với $ A $ qua gốc tọa độ $ O $. Tọa độ của $A'$ là $(-2;-5;-3)$}
	\loigiai{
		\begin{itemchoice}
			\itemch Ta có $A(2;5;3)$.
			\itemch Hình chiếu vuông góc của $A$ lên $Ox$ là $(2;0;0)$.
			\itemch Hình chiếu vuông góc của $A$ lên mặt phẳng $(Oyz)$ là $H(0;5;3)$.
			\itemch Điểm $ A' $ đối xứng với $ A $ qua gốc tọa độ $ O $ nên\\ $ \overrightarrow{OA'} =\overrightarrow{AO}\Rightarrow \overrightarrow{OA'}=-2\vec{i}-5\vec{j}-3\vec{k} \Rightarrow A'(-2;-5;-3)$.	
		\end{itemchoice}
	}
\end{ex}

\begin{ex}%[Câu 3]%[2H2H2-2]
\immini{Cho hình hộp chữ nhật $ OACB.O'A'C'B' $ có $ OA=2 $; $ OB=3 $; $ OC=4 $ được đặt trong không gian $ Oxyz $ như hình vẽ. Các khẳng định sau đúng hay sai?}
	{\begin{tikzpicture}[scale=0.6, font=\footnotesize, line join=round, line cap=round, >=stealth]
			\def\bc{4} % cạnh BC
			\def\ba{2} % cạnh BA
			\def\gocB{35} % góc B của đáy
			\coordinate[label=below right:$A$] (B) at (0,0);
			\coordinate (C) at (\gocB:\ba);
			\coordinate[label=below:$C$] (A) at (\bc,0);
			\coordinate[label=below right:$B$] (D) at ($(C)-(B)+(A)$);
			\coordinate[label=below left:$C'$] (A') at ($(A)+(90:\bc)$);
			\coordinate[label=left:$A'$] (B') at ($(B)-(A)+(A')$);
			\coordinate[label=above right:$O'$] (C') at ($(C)-(A)+(A')$);
			\coordinate[label=right:$B'$] (D') at ($(D)-(A)+(A')$);
			\draw (B')--(B)(D)--(D')--(A')--(B')--(C')--(D') (B)--(A)--(D) (A)--(A');
			\draw[dashed] (C)--(B)(C)--(D) (C)--(C');
			\foreach \diem in {A,B,C,D,A',B',C',D'}	\fill (\diem)circle(1.0pt);
			\draw[-stealth] (D)--+(0:1)node[right]{$y$};
			\draw[-stealth] (C')--+(90:1)node[above]{$z$};
			\draw[-stealth] (B)--+(-145:.75)node[below]{$x$};
			%\draw[-stealth,blue] (C)--+(0:.8)node[above]{$\vec{j}$};
			%\draw[-stealth,blue] (C)--+(90:.85)node[above left]{$\vec{k}$};
			%\draw[-stealth,blue] (C)--+(-145:.65)node[above]{$\vec{i}$};
			\fill(C)circle(1pt) node[below right]{$O$};
			\path (2,-1) node[below]{Hình 2.37};
\end{tikzpicture}}	
	\choiceTF
{\True Tọa độ điểm $A$ là $(2;0;0)$}
{\True Tọa độ điểm $C$ là $(2;3;0)$}
{\True Tọa độ của $C'$ là $(2;3;4)$}
{Gọi $ I $ là tâm của hình hộp chữ nhật đã cho. Tọa độ $ I $ là $H(1;2;2)$}
\loigiai{\begin{itemchoice}
		\itemch Ta có $A(2;0;0)$.
		\itemch Ta có $ C(2;3;0) $.
		\itemch  Ta có $ C'(2;3;4) $.
		\itemch $ I $ là trung điểm $ OC' $.\\
		$\overrightarrow{OI}=\dfrac{1}{2}\overrightarrow{OC'}=\dfrac{1}{2}\left(\overrightarrow{OA}+\overrightarrow{OB}+\overrightarrow{OO'}\right) =\dfrac{1}{2}(2\overrightarrow{i}+3\overrightarrow{j}+4\overrightarrow{k}) \Rightarrow I\left(1;\dfrac{3}{2};2\right)$.	
\end{itemchoice}
}
\end{ex}
\begin{ex}%[Câu 4]%[2H2V2-2]
\immini{Cho hình chóp $S.ABC$ có đáy $ABC$ là tam giác đều cạnh bằng $2$, $SA$ vuông góc với đáy và $SA =1$. Thiết lập hệ toạ độ $ Oxyz $ như hình vẽ bên. Các khẳng định sau đúng hay sai?}
{\begin{tikzpicture}[join = round, cap = round, thick, font = \footnotesize, scale =1]
		\path 
		(0:0) coordinate (A)
		++(0:4) coordinate (C)
		++(-160:3)coordinate(B)
		(A)++(90:2) coordinate (S)
		($(B)!1/2!(C)$) coordinate (O)
		($(S)+(O)-(A)$) coordinate (H)
		($(O)!1.4!(A)$) coordinate (y) node[above]{$y$}
		($(O)!1.4!(C)$) coordinate (x) node[above]{$x$}
		($(O)!1.4!(H)$) coordinate (z) node[right]{$z$}
		(intersection of S--C and O--H) coordinate (t)
		;
		\draw (S)--(A)--(B)--(C) (S)--(B) (O)--(z) (S)--(H) (C)--(t)
		;
		\draw[dashed] (C)--(A)--(O) (S)--(t)
		;
		\draw[->] (C)--(x); 
		\draw[->] (A)--(y);
		\draw[->] (H)--(z);
		\pic[draw,thin,angle radius=2mm] {right angle = S--A--O}
		pic[draw,thin,angle radius=2mm] {right angle = A--O--H}
		pic[draw,thin,angle radius=2mm] {right angle = O--H--S}
		pic[draw,thin,angle radius=2mm] {right angle = H--O--C}
		;
		\foreach \x/\g in {A/-100,C/-30,B/-90,O/-70,H/0,S/90}
		\fill (\x) circle (1pt)
		+(\g:3mm) node{$\x$};
\end{tikzpicture}}
\choiceTF
{\True Tọa độ điểm $ H $ là $\left(0;0;1\right)  $}
{Tọa độ điểm $ B $ là $ (1;0;0) $}
{\True Tọa độ điểm $ S $ là $ \left(0;\sqrt{3};1\right) $}
{Tọa độ véc-tơ $ \overrightarrow{AC}=\left(1;\sqrt{3};0\right) $}
\loigiai{
\begin{itemchoice}
\itemch Dựa vào hình vẽ ta thấy $ H \in Oz $, $ OH=AS=1 $ nên $ H(0;0;1) $.
\itemch  $ B $ thuộc tia đối của tia $ Ox $, $ OB=1 $ nên $ B\left(-1;0;0\right) $.
\itemch $ S \in \left(Oyz\right) $, hình chiếu của $ S $ lên trục $ Oy, Oz $ lần lượt là $ A, H $ và $ OA=\sqrt{AB^2-BO^2}=\sqrt{3} $; $ OH=AS=1 $ nên $ S\left(0;\sqrt{3};1\right) $.
\itemch  $ A(0;\sqrt{3};0) $, $ C(1;0;0) \Rightarrow \overrightarrow{AC}=\left(1;\sqrt{3};0\right)$.
\end{itemchoice}
}
\end{ex}

\Closesolutionfile{ans}
\inputansbox[3]{2}{ans/ans-2-B7-DS}
\caukq
\Opensolutionfile{ans}[ans/ans-2-B7-KQ]
\setcounter{ex}{0}
\begin{ex}%[Câu 1]%[2H2H2-2]
Trong không gian $Oxyz$, cho điểm $A(1;1;0)$, $B(-1;1;2)$, $C(3;1;4)$. Biết $ D(a;b;c) $ thỏa mãn $ABCD$ là hình bình hành. Tính $ a^2+b^2+c^2 $.
\shortans[0]{30}
\loigiai{ Gọi $D(x_D;y_D;z_D)$, ta có $ABCD$ là hình bình hành nên $\vec{AD}=\vec{BC}\Leftrightarrow (x_D-1;y_D-1;z_D)=(4;0;2)\Leftrightarrow D(5;1;2)$.\\
	Vậy $ a^2+b^2+c^2=30 $.
}
\end{ex}
\begin{ex}%[Câu 1]%[2H2H2-2]
	  Trong không gian $Oxyz$, cho hình lăng trụ tam giác $A B C.A' B' C'$ có $A(1 ; 0 ; 2)$, $B(3 ; 2 ; 5)$, $C'(11 ;-3 ; 8)$. Điểm $D'(a;b;c)$ sao cho $A B C D . A' B' C' D'$ là hình hộp. Tính $ a+b+c $.
	\shortans[0]{9}
	\loigiai{\immini{Vì  $A B CD . A' B' C'D'$ là hình hộp nên $ABC'D'$ là hình bình hành, suy ra $\overrightarrow{C'D'}=\overrightarrow{B A}$.\\
	Do đó $\heva{&x_D'-11=-2 \\& y_D'+3=-2 \\& z_D'-8=-3}$ hay $x_D'=9$, $y_D'=-5$ và $z_D'=5$.\\
	Suy ra $D'(9;-5;5)$.\\ 
	Vậy $ a+b+c=9 $.}
	{\begin{tikzpicture}[scale=0.6, font=\footnotesize, line join=round, line cap=round, >=stealth]
				\def\bc{3} % cạnh BC
				\def\ba{2} % cạnh BA
				\def\gocB{35} % góc B của đáy
				\coordinate[label=below left:$A'$] (A') at (0,0);
				\coordinate[label=above left:$B'$] (B') at (\gocB:\ba);
				\coordinate[label=below:$D'$] (D') at (\bc,0);
				\coordinate[label=right:$C'$] (C') at ($(B')-(A')+(D')$);
				\coordinate[label=above left:$B$] (B) at ($(B')+(90:\bc)$);
				\coordinate[label=left:$A$] (A) at ($(A')-(B')+(B)$);
				\coordinate[label=below right:$D$] (D) at ($(D')-(A')+(A)$);
				\coordinate[label=right:$C$] (C) at ($(B)-(B')+(C')$);
				\draw[dashed] (B')--(C');
				\draw[dashed] (B')--(B);
				\draw[dashed] (B')--(A');
				\draw(A')--(A)--(D)-- (A)--(B) (B)--(C)--(C')--(D')--(A') (D')--(D)--(C);
				\foreach \diem in {A,B,C,D,A',B',C',D'}	\fill (\diem)circle(1.0pt);
		\end{tikzpicture}}
	}
\end{ex}
\begin{ex}%[Câu 3]%[2H2V2-2]
	\immini{Ở một sân bay, vị trí của máy bay được xác định bởi điểm $M$ trong không gian $Oxyz$ như hình bên. Gọi $H$ là hình chiếu vuông góc của $M$ xuống mặt phẳng $(Oxy)$. Cho biết $OM = 50$, $\left(\overrightarrow{i},\overrightarrow{OH}\right) = 64^\circ$, $\left(\overrightarrow{OH},\overrightarrow{OM}\right) = 48^\circ$. Biết toạ độ của điểm $M$ là $ \left(a;b;c\right) $, $ a $, $ b $, $ c $ được làm tròn đến hàng phần chục. Tính $ a+b-c $.
	}{
	\begin{tikzpicture}[join = round, cap = round, thick, font = \footnotesize, scale =1]
		\path 
		(0,0) coordinate (O)++(0:4) coordinate (B)
		(O)++(90:2) coordinate (C)
		(O)++(-140:2) coordinate (A)		
		($(A)+(B)-(O)$) coordinate (H)
		($(H)+(C)-(O)$) coordinate (M)
		($(O)!1.3!(A)$) coordinate (x) node[below]{$x$}
		($(O)!1.2!(B)$) coordinate (y) node[below]{$y$}
		($(O)!1.3!(C)$) coordinate (z) node[right]{$z$}
		($(O)!0.3!(A)$) coordinate (i) node[below,scale =0.7]{$\overrightarrow{i}$}
		($(O)!0.3!(B)$) coordinate (j) node[below,scale =0.7]{$\overrightarrow{j}$}
		($(O)!0.3!(C)$) coordinate (k) node[right,scale =0.7]{$\overrightarrow{k}$}
		;
		\draw (O)--(A) (O)--(B) (O)--(C) (O)--(M);
		\draw[dashed] (C)--(M)--(H)--(O) (A)--(H)--(B)
		;
		\pic[draw,thin,angle radius=2mm] {right angle = M--C--O}
		pic[draw,thin,angle radius=2mm] {right angle = H--B--O}
		pic[draw,thin,angle radius=2mm] {right angle = H--A--O}
		;
		\draw[->] (A)--(x); 
		\draw[->] (B)--(y);
		\draw[->] (C)--(z);
		\draw[->,red] (O)--(i); 
		\draw[->,red] (O)--(j);
		\draw[->,red] (O)--(k);
		\foreach \x/\g in {A/130,H/-90,B/80,C/170,O/-90,M/45}
		\fill (\x) circle (1pt)
		+(\g:3mm) node[scale = 0.8]{$\x$};
	\end{tikzpicture}
	}
	\shortans[0]{7,6}
	\loigiai{
			Tam giác $OMH$ vuông tại $H$, $OM = 50$; $\widehat{MOH} = 48^\circ$ nên ta có $OH = OM\cdot \cos 48 \approx 33{,}5$ và $OC = MH = OM \cdot \sin 48 \approx 37{,}2$.\\
			Tam giác $OAH$ vuông tại $A$, $OH = 33{,}5$; $\widehat{AOH} = 64^\circ$ nên ta có $OA = OH\cdot \cos 64 \approx 14{,}7$,\\
			$OB = AH = OH\cdot \sin 64 \approx 30{,}1$.\\
			Suy ra 
			\begin{eqnarray*}
				\overrightarrow{OM} & = & \overrightarrow{OC} + \overrightarrow{OH} = \overrightarrow{OC} + \overrightarrow{OA}+\overrightarrow{OB} \\
				& = & 14{,}7\overrightarrow{i}+30{,}1\overrightarrow{j}+37{,}2\overrightarrow{k}.
			\end{eqnarray*}
		 $\Rightarrow M(14{,}7; 30{,}1; 37{,}2)$ và $ a+b-c=7{,}6 $.
	}
\end{ex}

\begin{ex}%[Câu 3]%[2H2V2-2]
	Trong không gian $Oxyz$, cho tam giác $ABC$ có $A(0;2;2)$, $C(4;-1;2)$; $ AB= \dfrac{15}{4}$, $ AC=5 $. Tìm cao độ của điểm $D$ là chân đường phân giác trong vẽ từ đỉnh $A$ của tam giác $ABC$.
	\shortans[0]{2}
	\loigiai{
		\immini{
			Theo tính chất chân đường phân giác ta có $\dfrac{DB}{DC}=\dfrac{AB}{AC}=\dfrac{3}{4}$, suy ra $$\overrightarrow{DB}=-\dfrac{3}{4}\overrightarrow{DC}\Leftrightarrow\heva{&x_B-x_D=-\dfrac{3}{4}\left(x_C-x_D\right)\\&y_B-y_D=-\dfrac{3}{4}\left(y_C-y_D\right)\\&z_B-z_D=-\dfrac{3}{4}\left(z_C-z_D\right)}\Leftrightarrow D(3;-1;2).$$
		}{	\begin{tikzpicture}[>=stealth,line join=round,line cap=round,line width=0.6pt,font=\footnotesize,scale=1]
			\path 
			(1,2)coordinate(A)node[above]{$A$}
			(0,0)coordinate(B)node[below]{$B$}
			(3,0)coordinate(C)node[below]{$C$}
			($(B)!.7!(C)$)coordinate(D)node[below]{$D$}
			%($(A)!.5!(C)$)coordinate(N)node[right]{$N$}
			%($(A)!.5!(B)$)coordinate(P)node[left]{$P$}
			;
			\draw (A)--(B)--(C)--cycle;
			\draw (A)--(D);
			\foreach \x in {A,B,C,D}\fill(\x)circle(1.5pt);
	\end{tikzpicture}}
	}
\end{ex}
\begin{ex}%[Câu 5]%[2H2C2-3]
	Trong không gian $Oxyz$, cho $A(1;0;0)$, $B(2;3;-1)$, $C(0;6;7)$ và gọi $M$ là điểm di động trên trục $Oy$. Tìm tung độ của điểm $M$ để $P = \left|\vec{MA}+\vec{MB}+\vec{MC}\right|$ đạt giá trị nhỏ nhất.
	\shortans[0]{$3$}
	\loigiai{
		$G$ là trọng tâm tam giác $ABC$, tọa độ $G(1;3;2)$.\\
		Gọi $M(0;y;0) \in Oy$, ta có
		$P=\left|\vec{MA}+\vec{MB}+\vec{MC}\right| = \left|3\vec{MG}\right| = 3MG$.\\
		Để $P_{\min} \Leftrightarrow MG_{\min} \Leftrightarrow M$ là hình của $G$ lên trục $Oy$. \\Khi đó $M$ có tọa độ $(0;3;0)$.
	}
\end{ex}
\begin{ex}%[Câu 6]%[2H2C2-2]
	Trong không gian $Oxyz$, cho điểm $A(0;-2;-3), B(-4;-4;1), C(2;-3;3)$. Giả sử điểm $M(a;0;c)$ trong mặt phẳng $Oxz$ thỏa mãn $P=MA^2+MB^2+2MC^2$ đạt giá trị nhỏ nhất. Tính $ a^2-c^2 $.
	\shortans[0]{$-1$}
	\loigiai{
Giả sử $ I $ thỏa mãn $ \overrightarrow{IA} +\overrightarrow{IB} +2\overrightarrow{IC} =\overrightarrow{0}$.\\
Khi đó tọa độ điểm $ I $ là $ \left(0;-3;1\right) $ và $ P=4MI^2+ IA^2+IB^2+2IC^2$.\\
Vì $ I $ cố định nên $ P $ đạt giá trị nhỏ nhất khi $ MI $ nhỏ nhất, hay $ M $ là hình chiếu của $ I $ trên $ (Oxz) $.\\
Vậy tọa độ điểm $M$ là $M(0; 0; 1).$
	}
\end{ex}
\Closesolutionfile{ans}
\inputansbox[3]{6}{ans/ans-2-B7-KQ}

